% Katalog B: die Box-Umgebungen (Palette Slot 2).
\ZSFNewColumn
\StartChapter[kat:boxen]{B · Boxen}

\begin{runintext}
  Zwölf Umgebungen, alle mit demselben Mustertext. Vier davon sind
  technisch ein \ZSFkeyword{defbox}-Preset und tragen ihre Äquivalenz in
  der Zweckzeile — dort sitzt die Frage, ob ein eigener Name nötig ist.
\end{runintext}

\SubsectionBar[kat:box-text]{Textbausteine}
\ZSFindex{Textbaustein}

\textVorBox{B-01 · runintext}
\begin{runintext}
  \Muster
\end{runintext}
\Zweck{Ruhiger Fliesstext zwischen strukturierten Blöcken, ohne jede Fläche.
Titellos, deshalb steht die ID davor.}

\begin{defbox}[B-02 · defbox]
  \Muster
\end{defbox}
\Zweck{Die allgemeine Inhaltsbox: Definition, Satz, gewichtige Aussage.}

\begin{warnbox}[B-03 · warnbox\ZSFTitleTag{Preset}]
  \Muster
\end{warnbox}
\Zweck{Stolperfalle mit Raum. Technisch defbox[tone=warn] plus vorbelegter
Titel »Achtung«.}

\begin{ZSFlist}[B-04 · ZSFlist]
  \ZSFItem{Mit Marke}{Mustertext zum Formvergleich.}
  \ZSFItem{Zweiter Eintrag}{Die zweite Zeile zeigt, wie der Container fasst.}
  \ZSFItem{Ohne Marke steht der Eintrag als blosse Aussage da.}
\end{ZSFlist}
\Zweck{Aufzählung von Fakten; die Marke jedes Eintrags ist optional.}

\textVorBox{B-05 · splitbox}
\begin{splitbox}[split=0.42]
  \MusterKurz
  \tcblower
  Die zweite Hälfte, getrennt durch tcblower.
\end{splitbox}
\Zweck{Zwei Blöcke nebeneinander, rahmenlos und titellos. Technisch defbox
mit fünf Reglern plus ruhiger Fläche.}

\SubsectionBar[kat:box-special]{Container mit eigenem Inhaltstyp}
\ZSFindex{Container}

\begin{runintext}
  Die folgenden vier tragen keinen freien Text, sondern einen bestimmten
  Inhaltstyp. Der Mustertext weicht deshalb ab — sie zeigen sonst nichts.
\end{runintext}

\begin{formulabox}[B-06 · formulabox]
  \[
    a^2 + b^2 = c^2
  \]
  \formulanote{Mustertext zum Formvergleich.}
\end{formulabox}
\Zweck{Formeln mit Kontext; betont geflächt in jedem Ton.}

\begin{tablebox}[B-07 · tablebox]
  \begin{ZSFtable}{Y{1.0} Y{1.0}}
    \ZSFheaderRow \ZSFhead{Erste} & \ZSFhead{Zweite} \\
    Mustertext & zum Formvergleich \\
    Die zweite Zeile & zeigt das Zebra
  \end{ZSFtable}
\end{tablebox}
\Zweck{Container für ZSFtable; läuft absichtlich bis an die Kante, damit das
Zebra nicht davor endet.}

\begin{figbox}[B-08 · figbox\ZSFTitleTag{Preset}]
  \Muster
\end{figbox}
\Zweck{Container für selbstgezeichnete Diagramme. Technisch
defbox[atomic]; hier ohne Bild, um den Container pur zu zeigen.}

\begin{goalbox}[B-09 · goalbox\ZSFTitleTag{Preset}]
  \GoalCondition{Mustertext zum Formvergleich}
  \GoalStep{\text{Zwischenschritt}}
  \GoalTarget{\text{Ziel}}
\end{goalbox}
\Zweck{Herleitung und Fallunterscheidung. Technisch defbox mit vier Reglern
plus ungetönter Fläche.}

\begin{valuegrid}[B-10 · valuegrid]{2}
  Mustertext & zum & Formvergleich \\
  Zweite Zeile & fixes & Gitter \\
  Erste Spalte & immer & Beschriftung
\end{valuegrid}
\Zweck{Kompakte Wertetabelle mit festem Gitter; das Argument ist die Anzahl
Datenspalten.}

\SubsectionBar[kat:box-nesting]{Boxen gruppieren und gliedern}
\ZSFindex{ZSFboxgroup}
\ZSFindex{ZSFinset}

\begin{ZSFboxgroup}[cols={Y{1.0} Q{1.0}},
                    head={\ZSFhead{Mustertext} & \ZSFhead{Formvergleich}}]
  \begin{tablebox}[B-11 · ZSFboxgroup]
    \begin{ZSFtable}[header=false]{\ZSFgroupcols}
      Erster Block & 1
    \end{ZSFtable}
  \end{tablebox}
  \begin{tablebox}[Zweites Mitglied]
    \begin{ZSFtable}[header=false]{\ZSFgroupcols}
      Zweiter Block & 2
    \end{ZSFtable}
  \end{tablebox}
\end{ZSFboxgroup}
\Zweck{Mehrere Boxen als ein Block: gemeinsame Kopfzeile, geschlossene
Zwischenräume, Spaltenaufteilung nur einmal gesetzt.}

\begin{tablebox}[B-12 · ZSFinset]
  \begin{ZSFtable}[header=false]{Y{1.0} Q{1.0}}
    Tabelle bis an die Kante & 1 \\
    damit das Zebra & 2
  \end{ZSFtable}
  \begin{ZSFinset}
    \Muster
  \end{ZSFinset}
\end{tablebox}
\Zweck{Rückt EINEN Block auf den normalen Innenabstand ein, wo die Box
absichtlich randlos ist.}
